% ============================================================
% 摘要
% 只写方法与结论的骨架；所有数字与具体结论用 \todo 占位。
% ============================================================

\begin{center}
  {\zihao{4}\bfseries 摘\quad 要}
\end{center}
\vspace{0.4em}

本文围绕商用啤酒风味代谢数据集，研究化合物谱与感官评分之间的映射关系、关键风味化合物的识别，以及给定风格下的化合物谱生成与跨风格迁移问题。针对附件数据样本量有限、化合物维度高、缺失值与异常值并存、感官评分为 0--6 离散等级等特点，本文采用“基线—改进—最终”的递进式建模路线，并在每一问中同时报告预测精度、稳定性与可解释性；同时引入题目建议的 Kaggle Top Beer Information 数据集，在 22 个风格层级做外部效度校验（覆盖 5,558 行中的 85.6\%），并明确其无法逐酒合并的限制。

针对问题一，本文以 \texttt{4.csv} 的 50 个感官维度为监督信号，把 \texttt{1.csv} 的 231 个化合物特征映射到连续感官潜变量，再校准为 0--6 分级。基线为逐维均值/岭回归；改进模型加入缺失指示、分位数变换与 \texttt{1.csv} 自带的风格先验；最终模型用 PLS2、Ridge、ExtraTrees、RandomForest 与 LightGBM 的 OOF 等权融合，并在训练折内用保序回归纠正预测收缩后做稳健 min-max 校准。重复 $5\times3$ 折交叉验证给出宏平均 $R^2=0.227$、Spearman $=0.432$；保序校准后的 QWK 为 0.376；30\% MCAR 下 $R^2$ 相对下降 19.3\%，MAR 下下降 10.9\%。

针对问题二，本文以 \texttt{2.csv} 构建 231 个化合物的共表达网络（Top-10 近邻，8 个社区，$Q=0.586$），用边际相关、弹性网、PLS、树模型与 LightGBM SHAP 多方法融合得到纯数据 Top-3，并用 bootstrap 入选频率（均值 0.525）与方法间 Jaccard（0.313）量化稳定性。由于纯数据 Top-3 在二乙酰、DMS、T2N 上选到的是风格/物化代理而非同名直接化合物，本文另给出机制导向 Top-3（文献先验 $\lambda=0.5$）并做 $\lambda$ 敏感性分析，把“预测性关联”与“化学机制”分开报告。

针对问题三，本文在 log1p 空间用条件高斯 copula 学习 22 个风格的化合物分布，并用高斯最优传输做跨风格迁移。最终模型对 DMS/Diacetyl/4VG/T2N 采用风格感知上限 $\max(\text{阈值},\text{风格 p75})$ 做投影；三风格平均 MMD 0.038、覆盖度 0.86，Pils/Lager $\to$ Tripel 的 Tripel 分类命中率从 0.00 提升到 0.91；严格阈值迁移下四类缺陷的超严格阈值率均为 0，生成模型在 style cap 下的四类缺陷超标率均为 0（严格阈值按风格取舍另行报告）；Tripel/Kriek/Lambic 的感官画像 Pearson 分别为 0.968/0.957/0.990。

最后，本文讨论了三问结果对啤酒风味设计与质量控制的启示：化合物模型适合做配方筛选与关键化合物监测，缺陷控制仍需前置到原料与发酵过程；外部数据只能做风格级校验，不能替代逐酒感官标签。

\vspace{0.8em}
\noindent\textbf{关键词：}啤酒风味；缺失数据；多任务回归；共表达网络；稳定性选择；条件生成模型；风味指纹
